\documentclass[track=junior,level=3,teacher]{kidscode}
\newcommand{\ostar}{\tikz[baseline=-0.6ex]\node[star,star points=5,star point ratio=2.3,draw=black,line width=0.8pt,fill=white,minimum size=0.6cm,inner sep=0pt]{};}
%% a script with a heading: \begin{refscript}{Sprite}{what it does} blocks \end{refscript}
\newenvironment{refscript}[2]{\par\needspace{8\baselineskip}\noindent\textbf{\color{kcaccent}#1:} #2\par\smallskip\begin{center}\begin{scratch}[scale=0.8]}{\end{scratch}\end{center}}
\begin{document}
\kctitle{Final Project Teacher Guide}{My Big Build: milestones, stars, the two reference projects and the graduation showcase}
\tableofcontents

%% =====================================================================
\section{The project in one page}
Each child chooses one of two paths and builds a project of their own, planned first on a
\textbf{plan card} (what it will look like, three parts in order, the must-haves, whose idea was
used) and checked by a partner. In week 10 each child shows it at the \textbf{graduation showcase},
explains one part of the code, and receives the junior graduation certificate.

\begin{center}\small\renewcommand{\arraystretch}{1.35}
\begin{tabularx}{\linewidth}{@{}>{\RaggedRight}X>{\RaggedRight\arraybackslash}X@{}}
\toprule
\textbf{Path A: Scratch game or art app} & \textbf{Path B: Python turtle postcard}\\
\midrule
Starts with the green flag. & Starts from the working starter \texttt{postcard.py} and runs without red text.\\
Uses a \textbf{repeat inside a repeat} or a \textbf{My Block}. & Uses a \textbf{\texttt{for} loop}.\\
Has a \textbf{score} or a \textbf{counter} (a variable). & Uses at least \textbf{two colors}.\\
 & Has \textbf{one shape} the child changed or added.\\
\midrule
\multicolumn{2}{@{}l@{}}{Both paths: planned on the plan card and checked by a partner against it.}\\
\bottomrule
\end{tabularx}
\end{center}

\begin{center}\small\renewcommand{\arraystretch}{1.35}
\begin{tabularx}{\linewidth}{@{}>{\bfseries\RaggedRight}p{3.0cm}>{\bfseries\RaggedRight}p{3.2cm}X@{}}
\toprule
When & Milestone & What the children do\\
\midrule
Week 9, warm-up (about 10 min) & M1 Plan card & Choose path A or B, draw the finished project, name three parts in order, check the must-haves, write whose idea they used.\\
Week 9, Make it (about 28 min) & M2 It works & Build one part at a time as Driver and Helper, test after every change; partner check against the plan card at about minute 42; save.\\
Week 10 (the whole session) & M3 Graduation showcase & Last bug check, a 3-minute showcase each, gallery walk with stickers, certificates.\\
\bottomrule
\end{tabularx}
\end{center}

\begin{materials}
  \item Print one \textbf{Project Brief} per child (the plan card is the page they fill in), plus a spare.
  \item Before week 9: build the Scratch reference (section \enquote{Reference A}, about 15 minutes) and
        run the Python reference (section \enquote{Reference B}) once on your Mac, so you can show both.
  \item The \textbf{core-path templates} for each child (see \enquote{Core path and stretch path} below):
        a Scratch \texttt{bigbuild-}nickname project, or a copy of \texttt{postcard.py} called
        \texttt{postcard\_}nickname\texttt{.py}. Not sure which path a child will choose? Prepare both; it takes a minute.
  \item Week 10: the week-10 show-me checklist, one sticker sheet per project (last pages), stickers,
        one graduation certificate per child (\texttt{extras/\allowbreak{}certificate}), and the progress trackers.
\end{materials}

\begin{teachernote}[Why this project]
The Big Build pulls every Level 3 goal together: patterns from loops and My Blocks (J3.G1), planning
in parts (J3.G2), a score or counter that counts something (J3.G3), reading and changing Python
(J3.G5), testing and debugging in blocks and text (J3.G6), and credit and privacy (J3.G7). Offering two
paths lets each child finish the junior track with something they are proud of: a child who loves
games keeps building in Scratch, a child who is curious about Python makes a first real program.
Both paths are equally valued at the showcase.
\end{teachernote}

%% =====================================================================
\section{Milestone walkthrough}

\subsection{M1 Plan card (week 9 warm-up, about 10 minutes)}
\begin{teacherscript}
\enquote{Today you start your Big Build, the last project of the junior track. You choose: a Scratch
game or art app, or a Python postcard. Let me show you one of each.} (Run the Scratch reference: click
the flag, press space a few times. Then run the Python reference in Thonny.) \enquote{Builders draw a
plan before they build a house. Coders do too. Our plan is called a plan card.}
\end{teacherscript}
\begin{enumerate}[itemsep=2pt]
  \item \textbf{Choose.} Read the two paths aloud. If a child cannot choose, ask: \enquote{Do you want
        people to play it, or to look at it?}
  \item \textbf{Draw} what it will look like when it is finished. Stick figures and a few words are perfect.
  \item \textbf{Name three parts} in the order they will build them. Scratch: \enquote{1. get ready,
        2. the flower pattern, 3. the counter.} Python: \enquote{1. the house, 2. the roof, 3. a sun.}
  \item \textbf{Check the must-haves} for their path, and write \textbf{whose idea} they used: the starter,
        a friend, a week of the course. Extra ideas go on the back of the card, for after M2.
\end{enumerate}
\begin{formativecheck}[M1 is done when]
A path is circled, there is a drawing, three parts are named in order, the must-haves for the path are
checked, and the credit line is filled in. Quick question: \enquote{Which part will you build first?}
\end{formativecheck}

\subsection{M2 It works (week 9 Make it, about 28 minutes)}
Each child builds their own project. The Helper holds the Driver's plan card, reads the next part
aloud, points at the category color in Scratch or the \texttt{CHANGE ME} line in Python, and says
\enquote{Test it!} after every change; swap on the timer. With one child, you are the Helper. Build in
this order and test after each step:
\begin{enumerate}[itemsep=2pt]
  \item \textbf{Open the template} (core path) or start fresh. Python: press \menu{Run} first, so the child
        sees that it already works.
  \item \textbf{Part 1, Part 2, Part 3,} one at a time. After each part: flag or Run, look, check it off on
        the plan card.
  \item \textbf{Partner check} (about minute 42): the partner runs the project and says \enquote{matches} or
        \enquote{not yet} for each part and must-have, then describes one \enquote{not yet} as a bug report:
        \enquote{I expected \dots, but it \dots} The owner fixes it, never the partner.
  \item \textbf{Save.} Scratch: \menu{File > Save to your computer}, in \texttt{Scratch Projects}. Python:
        \menu{Save} (the file already has the child's nickname).
\end{enumerate}
\begin{teacherscript}[Ask while you walk around]
\enquote{Which part of your plan card are you building now?} \enquote{Show me the code that makes this
part.} \enquote{What will change in the window when you press Run? Let's test it.}
\end{teacherscript}
\begin{formativecheck}[M2 is done when]
The project runs from a fresh start (flag or Run), has every must-have for its path, matches the three
parts more or less, the partner check is done with one fix made, and it is saved.
\end{formativecheck}

\subsection{M3 Graduation showcase (week 10, the whole session)}
\begin{enumerate}[itemsep=2pt]
  \item \textbf{Last bug check (10 minutes):} run the project twice, check the plan card, fix one thing
        only, save. Then the Python typo rescue from the lesson plan.
  \item \textbf{Practice (10 minutes):} you model a showcase in three sentences; children practice theirs.
  \item \textbf{Showcase (13 minutes):} each child presents for about 3 minutes while you fill in the
        week-10 show-me checklist.
  \item \textbf{Gallery walk and celebrate:} sticker sheets, stars, safety talk, graduation certificates.
\end{enumerate}
Most show-me rows appear during the bug check (typo fixing, running in Thonny) and the showcase
(plan card, explaining one part, credit).

%% =====================================================================
\section{Unblocking common problems}
\begin{center}\small\renewcommand{\arraystretch}{1.3}
\begin{longtable}{@{}>{\RaggedRight}p{5.2cm}>{\RaggedRight\arraybackslash}p{\dimexpr\linewidth-5.2cm-2\tabcolsep}@{}}
\toprule
\textbf{Problem} & \textbf{What to do}\\
\midrule\endhead
\bottomrule\endfoot
Cannot choose a path. & \enquote{Do you want people to play it, or to look at it?} Play: Scratch. Look: Python. Both are great.\\
The plan is huge. & Three parts is plenty. Circle the first three; the rest goes on the back for after M2.\\
\multicolumn{2}{@{}l}{\textbf{\color{kcaccent}Path A: Scratch}}\\ \midrule
The flower (or pattern) is one square only. & The square repeat sits under the outer repeat, not inside it (the week-1 bug). Drag it into the mouth of the outer repeat, above the turn.\\
The counter keeps counting from last time. & Add \enquote{set flowers to 0} under \enquote{when green flag clicked}.\\
The counter does not show on the stage. & In \menu{Variables}, check the box next to the variable.\\
A line joins all the flowers. & Pen up after the pattern, and pen down only after \enquote{go to mouse-pointer}.\\
A new script never runs. & It has no hat block on top, or it is under another sprite.\\
The My Block does nothing. & It is defined but never used: drag the pink \enquote{flower} block into the script.\\
The flowers are too big. & Make the cat smaller (\enquote{set size to 40 \%}) or the steps smaller (30 instead of 50).\\
\multicolumn{2}{@{}l}{\textbf{\color{kcaccent}Path B: Python}}\\ \midrule
Red text after a change. & Read the last red line together; it names the line. Undo the last change with \keys{Cmd}+\keys{Z} if needed, then make it again carefully. See the Teacher Guide troubleshooting table.\\
The roof floats above the house, or overlaps it. & Only one or two of the three \texttt{100}s were changed. Change all three to the same number.\\
\texttt{t.color(blue)} gives a \texttt{NameError}. & The quote marks are missing: \texttt{t.color("blue")}.\\
A new shape is drawn on top of the house. & Walk the turtle away first with the pen up: \texttt{t.penup()}, \texttt{t.forward(150)}, \texttt{t.pendown()}. Add one line at a time and Run.\\
A new shape is drawn once and the rest is a line. & The lines under \texttt{for} are not pushed in by the same 4 spaces.\\
The shape does not close. & The turn does not match the shape: use the turn card (square 90, triangle 120, hexagon 60).\\
\midrule
A real name is on the project or in the file name. & Change it together to a nickname. Explain why kindly.\\
\enquote{I want to add more!} & Great: first every must-have, the partner check and a save; then extras one at a time.\\
The project is lost. & Scratch: \menu{File > Load from your computer}. Thonny: \menu{Load}, look in \texttt{Coding}. Also check Downloads. Save every 10 minutes.\\
Tears or frustration. & Wiggle break. Then run it and find one thing that works: \enquote{Look, your house is already blue!}\\
\end{longtable}
\end{center}

%% =====================================================================
\section{The stars and how to use them kindly}
The brief has four rows, each with one to three stars. The stars are for the child, not a score to
report: there is no total, and children are never compared.

\begin{rubric}[My stars (adult version)]
\small\renewcommand{\arraystretch}{1.35}
\begin{tabularx}{\linewidth}{@{}>{\bfseries\RaggedRight}p{3.4cm}*{3}{>{\RaggedRight\arraybackslash}X}@{}}
 & \ostar & \ostar\,\ostar & \ostar\,\ostar\,\ostar\\
\midrule
It works & It starts, but parts do not work or there is red text. & Most parts work when the flag, a key or Run is used. & Every part works, every time, from a fresh start.\\ \midrule
It has every must-have & One must-have of the chosen path. & Most of them. & All of them (A: green flag, repeat in a repeat or My Block, score or counter; B: runs from the starter, \texttt{for} loop, two colors, one shape changed or added).\\ \midrule
I can explain one part of my code & Says what the project does overall. & Points to a script or loop and names some blocks or lines. & Picks one script or loop and says what each block or line does, in order.\\ \midrule
I planned it on my plan card & Filled in a plan card. & The project has some of the plan. & The project follows the three parts (changes along the way are fine).\\
\end{tabularx}
\end{rubric}

\begin{tip}[Using the stars kindly]
\begin{itemize}[itemsep=2pt]
  \item \textbf{The child colors first.} Read each row aloud and let them choose; then say what you saw.
        If you see more stars than they colored, say so. If you see fewer, ask \enquote{Can you show me?}
  \item \textbf{Two moments:} once at the end of week 9 (\enquote{what is left to do?}) and once after
        the showcase (\enquote{look how far you came}).
  \item \textbf{One star is a start.} Pick the one row to work on next and help with it right away.
  \item \textbf{Both paths count the same.} A three-part Scratch art app and a changed Python postcard
        can both earn every star.
\end{itemize}
\end{tip}

%% =====================================================================
\section{Reference A (Scratch): \enquote{Flower Garden}}
Build this yourself before week 9 so you can show it, then keep it on \textbf{your} Mac. It is one
complete answer for path A, and a longer version of Leo's Flower Garden from the workbook: children
make their own project, not a copy. It uses only whole numbers, a turn-card turn (60 for six squares
round a circle) and \enquote{go to mouse-pointer}: no x and y anywhere.

\subsection{Setting it up (no blocks yet)}
\begin{enumerate}[itemsep=2pt]
  \item New project (\menu{File > New}). The Cat is already there. Make it small: click the cat and type
        40 in the \textbf{Size} box below the stage.
  \item Click \menu{Add Extension} (bottom left) and choose \textbf{Pen}.
  \item Stage: click \menu{Choose a Backdrop} and pick a backdrop you like (a light one, so the flowers show).
  \item Variables (orange-red): \menu{Make a Variable}, name it \texttt{flowers}, \textbf{For all sprites}.
        Leave its box checked so the counter shows on the stage.
  \item Save as \texttt{Flower-Garden} in \texttt{Scratch Projects}.
\end{enumerate}

\subsection{The scripts, block by block (all on the Cat)}
\begin{refscript}{Script 1, get ready}{When the flag is clicked, the stage is cleared, the counter goes back to 0, the pen gets a color and a thick size (so everyone can see it), and the cat says what to do.}
\blockinit{when \greenflag clicked}
\blockpen{erase all}
\blockvariable{set \selectmenu{flowers} to \ovalnum{0}}
\blockpen{set pen color to \pencolor{kcblue}}
\blockpen{set pen size to \ovalnum{3}}
\blocklook{say \ovalnum{Press space for a flower!} for \ovalnum{2} seconds}
\end{refscript}

\begin{refscript}{Script 2, the My Block}{The pink My Block \enquote{flower} draws one flower: a square (repeat 4 of move and a quarter turn) inside a repeat 6, with a turn of 60 after each square. Six squares go once round, like the hexagon on the turn card.}
\blockmoreblocks{define \ovalmoreblocks{flower}}
\blockrepeat{repeat \ovalnum{6}}{
  \blockrepeat{repeat \ovalnum{4}}{
    \blockmove{move \ovalnum{30} steps}
    \blockmove{turn \turnright{} \ovalnum{90} degrees}
  }
  \blockmove{turn \turnright{} \ovalnum{60} degrees}
}
\end{refscript}

\begin{refscript}{Script 3, plant a flower}{When the space key is pressed, the cat jumps to the mouse-pointer with the pen up, draws one flower with the My Block, lifts the pen, adds 1 to the counter and pops. At 5 flowers it celebrates.}
\blockinit{when \selectmenu{space} key pressed}
\blockmove{go to \selectmenu{mouse-pointer}}
\blockpen{pen down}
\blockmoreblocks{flower}
\blockpen{pen up}
\blockvariable{change \selectmenu{flowers} by \ovalnum{1}}
\blocksound{start sound \selectmenu{Pop}}
\blockif{if \booloperator{\ovalvariable{flowers} = \ovalnum{5}} then}{
  \blocklook{say \ovalnum{A full garden!} for \ovalnum{2} seconds}
}
\end{refscript}

\begin{center}\small\renewcommand{\arraystretch}{1.3}
\begin{tabularx}{\linewidth}{@{}>{\bfseries}p{4.6cm}X@{}}
\toprule
Must-have (path A) & Where it is in the reference\\
\midrule
Green flag start & Script 1 gets everything ready.\\
A repeat inside a repeat, or a My Block & Both: the My Block \enquote{flower} holds a repeat 4 inside a repeat 6.\\
A score or a counter & The \texttt{flowers} variable: set to 0 in Script 1, changed by 1 in Script 3.\\
\bottomrule
\end{tabularx}
\end{center}

\begin{tip}[Questions to ask about Reference A]
\enquote{How do we make bigger flowers?} (Change 30 to 50.)
\enquote{What happens if the square repeat is under the repeat 6, not inside it?} (One square only: the week-1 bug.)
\enquote{Why does the counter need \enquote{set to 0}?} (Otherwise it keeps counting from last time.)
\enquote{Why pen up before the cat jumps?} (So no line joins the flowers.)
\enquote{Whose idea is the flower?} (Leo's, from the workbook, and the turn card: we say so at the showcase.)
\end{tip}

%% =====================================================================
\section{Reference B (Python): \enquote{A House in the Sun}}
This is one complete answer for path B. It starts from the starter \texttt{postcard.py} (an orange
house with a purple roof): the house and roof were made bigger (120), and the ground, a door, a sun
and three flowers were added one at a time, pressing Run after each. It uses only \texttt{forward},
\texttt{left}, \texttt{right}, \texttt{penup} and \texttt{pendown} to move, whole numbers, and turns from
the turn card. The flower is drawn by a \texttt{def}, the stretch path. Keep it on \textbf{your} Mac: it
is in \texttt{project/reference/}, which is never copied into the student folder. The file is tested
automatically (\texttt{test\_postcard\_reference.py}) and drawn in a real window before each release.

\begin{center}
%% alt: Sketch of the finished postcard: a long green ground line; on the left an orange square house with a blue door and a purple triangle roof; in the middle three small flowers with green stems and purple hexagon heads; top right a gold hexagon sun.
\begin{tikzpicture}[x=0.012cm,y=0.012cm,line cap=round]
  \draw[kcgreen,line width=2.4pt] (-250,0) -- (250,0);
  \draw[kcorange,line width=1.6pt] (-200,0) rectangle (-80,120);
  \draw[kcblue,line width=1.6pt] (-155,0) rectangle (-125,50);
  \draw[kcpurple,line width=1.6pt] (-200,120) -- (-80,120) -- (-140,224) -- cycle;
  \foreach \x in {-40,10,60} {
    \draw[kcgreen,line width=1.6pt] (\x,0) -- (\x,40);
    \draw[kcpurple,line width=1.6pt] (\x-10,40) -- ++(20,0) -- ++(60:20) -- ++(120:20) -- ++(180:20) -- ++(240:20) -- cycle;
  }
  \draw[kcyellow!80!black,line width=2.4pt] (130,180) -- ++(40,0) -- ++(60:40) -- ++(120:40) -- ++(180:40) -- ++(240:40) -- cycle;
\end{tikzpicture}
\end{center}

\pyfile{reference/postcard_reference.py}

\subsection{The program, part by part}
\begin{description}[style=nextline,leftmargin=1.2em,itemsep=3pt]
  \item[The docstring] The text in triple quotes at the top says what the program is for. Python
        does not run it; it is a note for people.
  \item[Part 1: the turtle] Exactly as in every week-6 to week-8 starter: \texttt{import turtle},
        make the turtle \texttt{t}, give it a shape and a speed.
  \item[Part 2: my own block (stretch)] \texttt{def flower():} gives a name to a group of lines, like a
        Scratch My Block. The lines inside are pushed in. It draws a green stem, steps 10 to the left so
        the petals sit in the middle, draws a purple hexagon (6 sides, turn 60 from the turn card), and
        walks back down the stem with the pen up, so it ends where it started. Defining a block does not
        draw anything yet: it draws when Part 8 uses it.
  \item[Part 3: the ground] The turtle starts in the middle, facing right. It walks 250 to the left with
        the pen up (a half turn, 180, then forward), turns back, and draws one long green line of 500.
  \item[Part 4: the house] The starter's square loop with bigger sides (120) and a thicker pen: 4 sides,
        a quarter turn (90) after each.
  \item[Part 5: the door] A rectangle is \enquote{long side, short side, two times}: a loop of 2 with two
        moves and two quarter turns. Stepping in 45 puts the door in the middle (45 + 30 + 45 = 120).
  \item[Part 6: the roof] The starter's triangle loop: 3 sides with the turn-card turn 120. Before it,
        the turtle walks back to the corner and up the wall with the pen up.
  \item[Part 7: the sun] A gold hexagon in the sky: the same loop as the flower petals, bigger.
  \item[Part 8: the flowers] The turtle walks down to the ground with the pen up. A \texttt{for} loop uses
        the block three times, with a step of 50 between the flowers.
  \item[The end] \texttt{t.hideturtle()} hides the turtle so the postcard looks finished, and
        \texttt{turtle.done()} keeps the window open until it is closed.
\end{description}

\begin{center}\small\renewcommand{\arraystretch}{1.3}
\begin{tabularx}{\linewidth}{@{}>{\bfseries}p{4.6cm}X@{}}
\toprule
Must-have (path B) & Where it is in the reference\\
\midrule
Starts from the starter and runs & The house (Part 4) and roof (Part 6) are the starter's loops.\\
A \texttt{for} loop & Every shape: house, door, roof, sun, petals, and the row of flowers.\\
At least two colors & Green, orange, blue, purple and gold.\\
One shape changed or added & Changed: the house and roof are bigger. Added: ground, door, sun, flowers.\\
Stretch: \texttt{def} & \texttt{flower()}, used three times.\\
\bottomrule
\end{tabularx}
\end{center}

\begin{tip}[Questions to ask about Reference B]
\enquote{Which loop draws the roof? How do you know?} (\texttt{range(3)} and the turn 120: a triangle.)
\enquote{How would you make five flowers?} (Change \texttt{range(3)} to \texttt{range(5)} in Part 8.)
\enquote{Why does the turtle lift its pen before walking?} (So no line is drawn between the shapes.)
\enquote{What would happen without \texttt{turtle.done()}?} (The window could close at once.)
\enquote{What did this postcard borrow?} (The house and roof from the starter: say so at the showcase.)
\end{tip}

%% =====================================================================
\section{Core path and stretch path}

\subsection{The core-path templates}
Make one per child before week 9 (about 3 minutes each).
\begin{itemize}[itemsep=2pt]
  \item \textbf{Path A (Scratch):} \menu{File > New}, add the Pen extension, pick a backdrop you like, keep
        the Cat, make a variable for the counter (\menu{Make a Variable}, For all sprites; the child can
        rename it later), and save as \texttt{bigbuild-} plus the child's nickname, for example
        \texttt{bigbuild-tiger}. The child adds the green-flag script, the pattern and the counter.
  \item \textbf{Path B (Python):} in Finder, open
        \texttt{Coding/\allowbreak kids-coding/\allowbreak grade-2-3/\allowbreak 3-bridge-to-python/\allowbreak project/\allowbreak starter/},
        duplicate \texttt{postcard.py} and rename the copy \texttt{postcard\_} plus the nickname, for example
        \texttt{postcard\_kai.py}. Never name a file \texttt{turtle.py}. The starter already runs: an orange
        house with a purple roof. Every spot to change is marked \texttt{\# CHANGE ME}; the child changes one,
        presses Run, and looks.
\end{itemize}

\subsection{The stretch path}
\begin{itemize}[itemsep=2pt]
  \item \textbf{Python:} a \texttt{def} for a shape the child draws more than once (a flower, a window, a
        tree), like the week-8 square block and Part 2 of Reference B. The brief shows a small \texttt{box()}
        example. Check: the lines inside \texttt{def} are pushed in, and the block is used by its name with
        brackets, \texttt{box()}.
  \item \textbf{Scratch:} a My Block for the repeated pattern (like \enquote{flower} in Reference A), or a My
        Block with a number input for the size.
\end{itemize}

%% =====================================================================
\section{Graduation showcase plan (week 10)}

\subsection{The graduation gallery}
The showcase is at home, for family only. Each project gets a \enquote{stand}: a Mac with the project open
(Scratch stage in full screen, or the Python program ready to Run in Thonny), the plan card beside it, a
sticker sheet, and a folded card with the project name and the maker's \textbf{nickname}. With 2--4
children, 10--15 minutes for the gallery walk is plenty.

\begin{center}\small\renewcommand{\arraystretch}{1.3}
\begin{tabularx}{\linewidth}{@{}>{\bfseries}p{3.0cm}X@{}}
\toprule
Role & Job\\
\midrule
Presenter (Driver) & Runs the project, points with the mouse and says the three showcase sentences.\\
Partner (Helper) & Holds up the presenter's plan card while they explain, and starts the clapping. Swap for the next project.\\
Visitors & Ask \enquote{May I try it?}, run it, put one sticker in the column they liked best, and say why. Only kind words; visitors never change someone else's project.\\
Host (you) & Welcomes the family, keeps time, fills in the show-me checklist, makes sure every child gets the same attention, reads the safety reminder.\\
\bottomrule
\end{tabularx}
\end{center}

\subsection{The three showcase sentences}
Write them on a sticky note for each child. Practice once before the family arrives. Children who do not
read yet can point to pictures you draw next to each line.
\begin{enumerate}[itemsep=2pt]
  \item \textbf{My plan:} \enquote{My plan card said \underline{\hspace{4cm}}. \underline{\hspace{1.2cm}} parts work.}
  \item \textbf{One part of my code:} point at one script or loop: \enquote{This \underline{\hspace{2.5cm}} makes
        \underline{\hspace{3cm}}. First it \underline{\hspace{2.5cm}}, then it \underline{\hspace{2.5cm}}.}
  \item \textbf{What I changed, and credit:} \enquote{I changed \underline{\hspace{3cm}}. The
        \underline{\hspace{2.5cm}} came from \underline{\hspace{3cm}}.}
\end{enumerate}
Then: \enquote{Do you have a question for me?}

\subsection{Certificates and closing}
Read the sticker sheets aloud to each child, then hand out the graduation certificate (write the nickname
or first name by hand). Tell them what comes next: main-track Level 1, where Scratch, loops and Python will
feel familiar. Close with a round of \enquote{one thing I am proud of, one thing I want to code next},
color the last stars on the progress tracker, and finish with a big clap.

\begin{safety}[Showcase safety]
\begin{itemize}[itemsep=1pt]
  \item Family only, at home, on the screen. Nothing is uploaded, posted or emailed.
  \item Nicknames only on the screen, in file names and on the stand cards.
  \item Photos only if the child agrees, with no names on the screen, and kept in the family.
  \item Credit: every project says what it borrowed (the starter, a friend's idea, the Scratch library).
  \item Next steps safely: if a website or app ever asks for your name, age, school or a photo, stop and
        ask a grown-up.
\end{itemize}
\end{safety}

\begin{teachernote}[Fitting it into the 60-minute session]
0--10 last bug check and the typo rescue (warm-up); 10--20 you model the three sentences, children
practice (show me); 20--33 showcases, week-10 checklist (make it); 33--35 wiggle break; 35--50 gallery
walk with sticker sheets, show-me catch-up (make it); 50--60 safety talk, what is next, certificates (share).
\end{teachernote}

\clearpage
\section{Sticker sheets (print and cut)}
One per project. Visitors put a sticker in the column they liked best, and a grown-up can write one more idea.
\newcommand{\stickersheet}{%
%% alt: A sticker sheet that says Your Big Build! with three columns of dashed sticker spots under the colors, the pattern and the surprise, and a line for one more idea.
  \begin{tikzpicture}[spot/.style={circle,draw=kcgray,dashed,minimum size=0.8cm,inner sep=0pt}]
    \draw[kcpurple,line width=1.2pt,rounded corners=6pt] (0,0) rectangle (16.5,8.2);
    \node[anchor=north west,font=\bfseries] at (0.3,8.0) {\faGraduationCap\ Your Big Build!};
    \node[anchor=north west,font=\small] at (0.3,7.3) {Put a sticker in the column you liked best.};
    \foreach \icon/\word [count=\i] in {\faPalette/the colors,\faTh/the pattern,\faGift/the surprise}
      {\node[font=\bfseries] at (\i*5.5-2.75,6.2) {\icon\ \word};
       \draw[kcgray!60] (\i*5.5-5.3,6.7) rectangle (\i*5.5-0.2,1.6);
       \foreach \r in {1,2,3}
         \foreach \c in {1,2,3}
           \node[spot] at (\i*5.5-5.5+\c*1.3,6.2-\r*1.3) {};}
    \node[anchor=north west,font=\small] at (0.3,1.3) {One more idea:};
    \draw[kcgray!60] (3.0,0.5) -- (16.0,0.5);
  \end{tikzpicture}}
\begin{center}
\stickersheet\par\bigskip
\stickersheet
\end{center}
\end{document}
